\section{Why the degree is exactly \texorpdfstring{\(h\)}{h}}\label{sec:exact-degree}
\source{Fauzan \cite[(2.9)]{fauzan}; Lai \cite[Proposition 6.1]{lai};
the same formula is stated in Persiflage \cite{persiflage}.}

The main proof only needs \(\deg\Delta_K\le h\). The residue pairing
also shows why the highest term survives: there are exactly \(h\)
remaining poles, matching the dimension of the polynomial space.
Put \(a_k=-k^2\) and \(V_{ik}=a_k^i\), with \(0\le i<h\) and \(N<k\le K\).
By \eqref{eq:gram-residue}, if \(G_K(X)=G_0+XG_1\), then
\[
 G_1=V\operatorname{diag}(d_k)V^{\mathsf T},\qquad
 d_k=\frac{k^4D_N(a_k)^5}{T_K'(a_k)}.
\]
Every \(d_k\) is nonzero. The nodes \(a_k\) are distinct, so evaluation
at them is injective on polynomials of degree less than \(h\):
a nonzero such polynomial cannot have \(h\) roots. Thus \(V\) is
invertible and \(\det G_1\ne0\). Taking the \(X\) term from every factor
in the determinant expansion gives
\[
 [X^h]\Delta_K(X)=\det G_1\ne0,\qquad \deg\Delta_K=h.
\]
For the explicit leading coefficient, the Vandermonde formula and
pairing the factors in the derivatives give
\[
 \det V=\prod_{k<\ell}(a_\ell-a_k),\qquad
 \prod_k T_K'(a_k)=(-1)^{h(h-1)/2}(\det V)^2.
\]
Hence
\[
 [X^h]\Delta_K(X)
 =(-1)^{h(h-1)/2}\prod_{k=N+1}^K k^4D_N(-k^2)^5
 \in\Z\setminus\{0\}.
\]
This concerns one coefficient; the arithmetic must still control them all.
